% Discussion and Conclusion

This chapter interprets the experimental findings from RQ1, discusses limitations of the current work, and outlines future research directions including RQ2 and machine learning extensions.

\section{Interpretation of Results}

\subsection{Communication Range Effects on Harvesting Efficiency}

The experimental results provide clear evidence that communication range has a dramatic effect on harvesting performance. The key findings are:

\begin{itemize}
    \item \textbf{Communication is essential:} Agents with no communication (range 0) harvest only 354.61 fruits on average, compared to 487-500 fruits with communication. This represents a 37-41\% improvement.

    \item \textbf{Optimal range is range 2:} While ranges 4, 6, and 8 achieve slightly higher yields (499-500 fruits), range 2 is most efficient, achieving 97.6\% of maximum yield while using only 67.5\% of the messages required by range 8.

    \item \textbf{Diminishing returns at larger ranges:} Ranges 4, 6, and 8 show a plateau effect with nearly identical yields (within 1.2 fruits) despite increasingly higher message overhead. This indicates that beyond range 4, additional communication provides minimal benefit.
\end{itemize}

These results directly answer RQ1: communication range significantly affects harvesting efficiency, with an optimal range that balances coordination benefits against communication costs.

\subsection{Comparison with Literature}

These findings align with research on swarm robotics and multi-agent coordination. Similar to ant colonies that use pheromone trails (a form of local communication), our agents benefit from sharing information within a limited range. The sweet spot we observe (range 2) is consistent with the principle that local communication is often more efficient than global broadcast, as discussed in the literature on distributed systems and swarm intelligence.

The diminishing returns at larger ranges support the theoretical understanding that communication overhead can outweigh coordination benefits. This is particularly relevant for resource-constrained systems where message passing has computational costs.

\subsection{Robustness Across Configurations}

The experimental results show that the communication range effect is robust and consistent:

\begin{itemize}
    \item The pattern holds across different team sizes (10 and 20 agents)
    \item The pattern holds across different resource densities (0.15 and 0.25)
    \item The relative ranking of ranges remains consistent: 0 < 2 < 4 ≈ 6 ≈ 8
\end{itemize}

This consistency suggests that the findings are not artifacts of specific parameter choices but reflect fundamental principles of multi-agent coordination.

\section{Limitations}

This study has several limitations that should be acknowledged:

\begin{itemize}
    \item \textbf{Simplified agent behaviors:} Agents use greedy movement and random walk, not sophisticated planning or learning algorithms.
    \item \textbf{Grid-based abstraction:} The discrete grid environment is simpler than continuous space with realistic physics.
    \item \textbf{Homogeneous agents:} All agents are identical with the same capabilities and behaviors.
    \item \textbf{Static environment:} No obstacles, no dynamic changes beyond fruit harvesting.
    \item \textbf{Static resources only:} Experiments focus on non-regenerating fruit; regenerating resources are not tested.
    \item \textbf{Deterministic communication:} All agents within range receive messages; no communication failures or noise.
\end{itemize}

\section{Future Work}

Several promising directions emerge from this research:

\subsection{Resource Dynamics (RQ2)}

The original RQ2 asks whether resource dynamics affect optimal communication range. Future work should implement and test Replenishing mode where fruit regenerates with some probability. We hypothesize that sustained communication becomes more valuable in replenishing environments, potentially shifting the optimal range higher.

\subsection{Machine Learning Extensions}

\begin{itemize}
    \item \textbf{Reinforcement Learning:} Train agents to learn when and what to communicate using Q-learning or policy gradient methods, potentially discovering adaptive communication strategies.
    \item \textbf{Deep Learning:} Use neural networks to predict optimal communication ranges from environment properties (density, team size, grid size).
    \item \textbf{Multi-Agent Reinforcement Learning (MARL):} Investigate how agents can learn communication policies in a decentralized manner.
    \item \textbf{Transfer Learning:} Train on one environment and transfer to different team sizes or resource densities.
\end{itemize}

\subsection{Realism and Validation}

\begin{itemize}
    \item Extend to continuous space with realistic movement dynamics.
    \item Add communication constraints (bandwidth limits, message delays, noise).
    \item Implement heterogeneous agents with different capabilities.
    \item Validate findings with physical robot experiments.
\end{itemize}

\section{Conclusion}

This dissertation investigated RQ1: "Does communication range affect harvesting efficiency in multi-agent systems?" Through systematic experimentation with 400 simulation runs across five communication ranges, two team sizes, and two resource densities, we found that communication range significantly affects harvesting performance. The results demonstrate that there is an optimal communication distance that balances the benefits of information sharing against the costs of message overhead.

These findings contribute to our understanding of multi-agent coordination and provide a foundation for future work on adaptive communication strategies, resource dynamics, and machine learning approaches to agent communication.

